%-------------------------------------------------------%
%                         英文摘要
%-------------------------------------------------------%
\chapter*{\textbf{ABSTRACT}}
\eabstractpagestyle % 中文摘要的页眉页脚
Under the rapid development of big data, deep learning, artificial intelligence, 
and intelligent network technologies, this study proposes a dynamic holding control 
strategy based on deep reinforcement learning to address common operational efficiency 
issues in urban bus systems such as bus bunching and unstable headway. Aiming at the
shortcomings of traditional holding control methods in real-time response and dynamic 
environment adaptation, we innovatively construct a multi-agent collaborative optimization 
framework that integrates intelligent connected vehicle data with reinforcement learning algorithms 
to achieve real-time dynamic holding control for bus operations.

By analyzing the actual operating environment of buses, we abstract and model core elements within it. 
Simultaneously, we conduct comprehensive elaboration and detailed analysis of reinforcement learning 
algorithms, and the algorithm framework Actor-Critic suitable for this research problem is selected. 
We adopt the Multi-Agent Proximal Policy Optimization (MAPPO) as the base algorithm. 
Based on this, we build a bus fleet operation model and meticulously design a bus operation simulation 
experimental process. Through investigating the causes of headway instability in bus operations, 
we propose a holding control strategy aimed at stabilizing headway.

This paper models bus holding control as a partially observable Markov game.
By designing an Event Graph Network (EGN) to describe the asynchronous timing characteristics 
of vehicle arrivals, the Graph Attention Network (GAT) is used to learn the event graph 
structure information and quantify the influence intensity of spatio-temporal interaction 
between vehicles, to supplement the limitation of the Markov game that the influence of 
other agents' actions cannot be considered due to the local observation independence constraint.
Building upon this framework, we propose an enhanced Multi-Agent Proximal Policy Optimization (MAPPO) 
algorithm by integrating a graph attention network (GAT) into the original architecture, 
thereby establishing a dynamic holding control model with environmental awareness and collaborative 
decision-making capabilities.
In this model, each bus is regarded as an agent, and the multi-dimensional state space 
(including vehicle operational stability metrics and passenger demand characteristics)
and continuous action space (holding duration decisions) are designed, and a composite reward function is 
constructed that integrates the headway stability and the minimization of passenger waiting time.

To validate the proposed model's performance, a one-way bus circular bus route was designed and implemented. 
Comparative experiments were conducted against baseline models using predefined evaluation metrics, 
demonstrating the superiority of the proposed approach. To further verify its practical feasibility, 
the model was extended to a real-world conventional bus line scenario incorporating additional operational complexities. 
Using the Chengdu T34 bus route as a validation platform, the model was compared with rule-based control, 
mathematical optimization control, and the original MAPPO algorithm across both micro and macro perspectives.
The numerical results indicate that the bus holding control model proposed in this paper, 
based on the improved MAPPO algorithm, shows performance improvements over the No-holding control (NH), 
the naive hard-holding control (HH), the forward headway-based holding control (FH), 
the backward headway-based holding control (BH), and the control model based on the original MAPPO algorithm, 
with respect to the average waiting time (AWT) being improved by 34.95\%, 27.97\%, 12.56\%, 20.00\%, and 7.84\%, 
respectively. For the average travel time (ATT) metric, the improvements are 4.3\%, 21.25\%, 16.63\%, 8.99\%, and 2.94\%, 
respectively. The average occupancy dispersion (AOD) is reduced by 
77.14\%, 5.88\%, 11.11\%, 20.00\%, and 11.11\%, respectively (lower values indicate better load balancing).
Regarding the average holding time (AHT), the performance compared to 
the naive hard-holding control (HH), the forward headway-based holding control (FH), 
the backward headway-based holding control(BH), and the control model based on the original MAPPO algorithm 
increased by 74.51\%, 69.05\%, 45.83\%, and 27.78\%, respectively. 
The results suggest that the proposed model performs well in maintaining the stability of bus headways, 
reducing bus bunching, lowering the cost of passenger travel time, improving the utilization of 
public transport resources, and enhancing the quality of public transport services.



% Numerical results revealed that the improved MAPPO-based bus holding control model achieved significant improvements 
% compared to No-holding control (NH), Naive hard-holding control (HH), Forward headway-based holding control (FH), 
% Backward headway-based holding control(BH), and original MAPPO-based holding control:
% \begin{itemize}
%     \item Average Waiting Time (AWT): 34.95\%, 27.97\%, 12.56\%, 20.00\%, and 7.84\% reduction respectively
%     \item Average Travel Time (ATT): 4.3\%, 21.25\%, 16.63\%, 8.99\%, and 2.94\% reduction respectively
%     \item Average Occupancy Dispersion (AOD): 77.14\%, 5.88\%, 11.11\% 20.00\%, and 11.11\% improvement respectively
%     \item Average Holding Time (AHT): 74.51\%, 69.05\%, 45.83\%, and 27.78\% reduction respectively
% \end{itemize}
% These findings demonstrate that the proposed model effectively maintains headway stability, 
% mitigates bus bunching incidents, reduces passenger travel time costs, optimizes resource utilization, 
% and enhances overall service quality in public transit systems.


\vspace{1em}

\noindent \hangindent=4.9em \textbf{Keywords: }Reinforcement learning, Bus bunching, Agent, Holding control, Asynchronous characteristics